% ============================================================
% 05-eksperimentalna-analiza.tex — PRVA VERZIJA
% ============================================================

Eksperimentalna analiza ima za cilj usporediti u\v{c}inkovitost
jedanaest SMOTE varijanti na skupovima podataka razli\v{c}itih
karakteristika. Analizirat \'{c}e se utjecaj omjera neuravnote\v{z}enosti,
dimenzionalnosti i prisutnosti \v{s}uma na relativne performanse
algoritama, a statisti\v{c}kim testovima provjerit \'{c}e se
zna\v{c}ajnost uo\v{c}enih razlika.

\subsection{Opis skupova podataka}

Opisuju se skupovi podataka kori\v{s}teni u eksperimentalnoj
analizi. Koristit \'{c}e se kombinacija stvarnih i sinteti\v{c}kih
skupova. Za stvarne skupove bit \'{c}e odabran odre\dj{}eni broj
neuravnote\v{z}enih skupova podataka iz razli\v{c}itih domena, pri
\v{c}emu \'{c}e se nastojati obuhvatiti \v{s}to raznovrsniji raspon
omjera neuravnote\v{z}enosti, dimenzionalnosti i veli\v{c}ina kako bi
se dobio \v{s}to op\'{c}enitiji uvid u pona\v{s}anje algoritama.
Tako\dj{}er, poku\v{s}at \'{c}e se identificirati skupovi podataka koji
donekle prate odre\dj{}ene statisti\v{c}ke distribucije (primjerice,
normalnu, eksponencijalnu ili multimodalnu) kako bi se ispitao utjecaj
samih algoritama na podatke razli\v{c}itih distribucijskih svojstava.

Sinteti\v{c}ki skupovi generirat \'{c}e se s kontroliranim parametrima --- omjerom neuravnote\v{z}enosti, dimenzionalno\v{s}\'{c}u i razinom
\v{s}uma --- \v{c}ime \'{c}e se omogu\'{c}iti ciljano ispitivanje
utjecaja pojedinih karakteristika na performanse algoritama, neovisno
o specifi\v{c}nostima stvarnih domena. Osvrnuti se na ograni\v{c}enja
sinteti\v{c}kih skupova podataka.

Dat \'{c}e se tablica s karakteristikama svih kori\v{s}tenih skupova:
naziv, broj primjera, broj zna\v{c}ajki, omjer neuravnote\v{z}enosti
te kategorija prema IR-u (nizak, srednji, visok).

\subsection{Postavke eksperimenta}

Eksperimenti \'{c}e se provesti kori\v{s}tenjem 5-struke stratificirane
unakrsne validacije, pri \v{c}emu je u svakoj iteraciji omjer
neuravnote\v{z}enosti o\v{c}uvan u skupu za treniranje i testiranje. Za
svaki SMOTE algoritam testirat \'{c}e se vrijednosti parametra
$k \in \{3, 5, 7, 10\}$, a preuzorkovanje \'{c}e se vr\v{s}iti do
potpune balansiranosti klasa. Klasifikatori uklju\v{c}uju stablo odluke,
slu\v{c}ajnu \v{s}umu sa 100 stabala, XGBoost, logisti\v{c}ku regresiju,
SVM s RBF jezgrom, $k$-NN s $k=5$, Gaussov naivni Bayes i vi\v{s}eslojni
perceptron (MLP). XGBoost koristi parametar \texttt{scale\_pos\_weight}
za nativno rukovanje neuravnote\v{z}eno\v{s}\'{c}u klasa, dok MLP
s jednim skrivenim slojem od 100 neurona slu\v{z}i kao neuronska
referentna to\v{c}ka. Kao bazne metode koristit
\'{c}e se klasifikacija bez preuzorkovanja, nasumi\v{c}no poduzorkovanje
i nasumi\v{c}no preuzorkovanje.

Prije klasifikacije provest \'{c}e se standardizacija zna\v{c}ajki
pomo\'{c}u \texttt{StandardScaler}-a. Svaka konfiguracija ponovit
\'{c}e se 30 puta s razli\v{c}itim po\v{c}etnim vrijednostima generatora
slu\v{c}ajnih brojeva. Prikupljat \'{c}e se F1, G-Mean, AUC-ROC,
Balanced Accuracy i MCC metrike, a rezultati \'{c}e se prikazati kao
aritmeti\v{c}ka sredina $\pm$ standardna devijacija kroz svih 30
ponavljanja. Navest \'{c}e se i hardverska te softverska konfiguracija
na kojoj su eksperimenti izvedeni.

\subsection{Rezultati}

Zbog velikog broja kombinacija algoritama, klasifikatora i skupova
podataka, rezultati \'{c}e se prikazati na vi\v{s}e komplementarnih
na\v{c}ina, svaki s razli\v{c}itom razinom granularnosti. U glavnom
tekstu fokus \'{c}e biti na agregiranim i statisti\v{c}ki obra\dj{}enim
prikazima koji omogu\'{c}uju uo\v{c}avanje \v{s}irih obrazaca, dok
\'{c}e kompletni rezultati svih eksperimenata biti dani u prilogu.

Prvo, za svaku metriku dat \'{c}e se zbirna tablica s prosje\v{c}nim
rangom svakog algoritma kroz sve skupove podataka i klasifikatore,
poredana od najboljeg prema najlo\v{s}ijem. Ovakav prikaz omogu\'{c}uje
brzi uvid u ukupni poredak algoritama bez zatrpavanja detaljima.
Zatim \'{c}e se za svaku metriku prikazati kriti\v{c}ni dijagram
(CD dijagram) koji vizualno sa\v{z}ima rezultate Friedmanovog i Nemenyi
testa --- algoritmi koji se statisti\v{c}ki zna\v{c}ajno ne razlikuju
povezani su vodoravnom crtom.

Za detaljniji uvid, prikazat \'{c}e se heatmap-ovi koji za svaki par
(algoritam, skup podataka) prikazuju vrijednost odabrane metrike uz
kodiranje bojom, \v{s}to omogu\'{c}uje uo\v{c}avanje obrazaca poput
toga koji algoritmi dominiraju na skupovima s visokim IR-om, a koji na
onima s niskim. Box-plotovi \'{c}e se koristiti za prikaz distribucije
rezultata po algoritmima kroz sva ponavljanja, \v{c}ime se dobiva uvid
ne samo u prosje\v{c}nu performansu ve\'{c} i u njenu varijabilnost.

Za svaki klasifikator izdvojit \'{c}e se i kratko prokomentirati
najva\v{z}niji uo\v{c}eni obrasci --- primjerice, kod kojih klasifikatora
preuzorkovanje donosi najve\'{c}e pobolj\v{s}anje, kako se pona\v{s}aju
geometrijska pro\v{s}irenja u odnosu na klasi\v{c}ne varijante te u kojim
situacijama SMOTE nema u\v{c}inka ili \v{c}ak pogor\v{s}ava rezultate.

Kompletne tablice sa svim vrijednostima metrika za svaku kombinaciju
algoritma, klasifikatora i skupa podataka bit \'{c}e prilo\v{z}ene u
prilogu, organizirane po metrikama radi preglednosti.


\subsection{Statisti\v{c}ka analiza}

Za statisti\v{c}ku usporedbu performansi algoritama koristit \'{c}e se
metodologija Dem\v{s}ara \cite{demsar2006}. Provest \'{c}e se
Friedmanov test za utvr\dj{}ivanje postoje li zna\v{c}ajne razlike
me\dj{}u algoritmima, Nemenyi post-hoc test za identifikaciju razlika
izme\dj{}u konkretnih parova (prikazano kriti\v{c}nim dijagramima) te
Wilcoxon signed-rank test za usporedbu svake izvedenice s osnovnim
SMOTE-om. Dodatno, skupovi \'{c}e se grupirati prema omjeru
neuravnote\v{z}enosti, dimenzionalnosti i razini \v{s}uma kako bi se
ispitalo kako ovi faktori utje\v{c}u na relativne performanse
algoritama.

\subsection{Diskusija}

Na temelju provedene eksperimentalne i statisti\v{c}ke analize
interpretirat \'{c}e se dobiveni rezultati s naglaskom na sljede\'{c}a
pitanja: koji SMOTE algoritam pokazuje najbolje ukupne performanse,
u kojim situacijama se preporu\v{c}uje koja izvedenica, kako
geometrijska pro\v{s}irenja stoje u odnosu na klasi\v{c}ne varijante,
kakav je utjecaj parametra $k$ na razli\v{c}ite klasifikatore te koja
je ra\v{c}unska zahtjevnost pojedinih algoritama. Dani \'{c}e se i
prakti\v{c}ne preporuke za odabir algoritma temeljem karakteristika
skupa podataka. Rezultati \'{c}e se usporediti s nalazima iz literature
\cite{bajer2019,dudjak2020,kovacs2019,haixiang2017}.
